\chapter{Funzioni di più variabili, covarianza e correlazione}

Nel capitolo precedente abbiamo descritto una coppia di variabili aleatorie con la sua pmf congiunta. Ora le usiamo:
combiniamo $X$ e $Y$ in una nuova variabile $Z = g(X, Y)$ (una somma, un prodotto, una formula qualsiasi) e ne
calcoliamo pmf e media; scopriamo che la media di una somma è sempre la somma delle medie; e introduciamo due numeri,
la \textbf{covarianza} e il \textbf{coefficiente di correlazione}, che riassumono in modo semplice quanto due
variabili «si muovono insieme».

% ============================================================
\section{Funzioni di una coppia di variabili}
% ============================================================

Sia $(X, Y) \sim p_{X,Y}(x, y)$, $(x, y) \in \mathcal{X} \times \mathcal{Y}$, una variabile doppia, e sia $g(x, y)$ una
funzione di due variabili il cui dominio include $\mathcal{X} \times \mathcal{Y}$. Si vuole caratterizzare
$Z = g(X, Y)$ in termini di pmf e di media statistica. Per prima cosa si determina l'alfabeto $\mathcal{Z}$, l'insieme
dei valori $g(x, y)$. Poi si ragiona esattamente come nel Capitolo~6 per le funzioni di una variabile.
\begin{itemize}
  \item Se $\card{\mathcal{Z}} = \card{\mathcal{X}}\,\card{\mathcal{Y}}$, la funzione è biunivoca sulle coppie: per ogni
        $z$ esiste un unico punto $(x(z), y(z))$ con $z = g(x, y)$, e
        \[ P(Z = z) = p_Z(z) = p_{X,Y}\left[ x(z), y(z) \right] . \]
  \item Se $\card{\mathcal{Z}} < \card{\mathcal{X}}\,\card{\mathcal{Y}}$, la relazione precedente vale nei punti in cui
        l'inversa è unica, mentre negli altri le probabilità «collassano»: se $z = g(x, y)$ per tutte le coppie di un
        insieme $A(z) \subseteq \mathcal{X} \times \mathcal{Y}$, allora
        \[ p_Z(z) = \sum_{(x, y) \in A(z)} p_{X,Y}(x, y) . \]
\end{itemize}

\begin{formula}[pmf di una funzione di due variabili]
\[ p_Z(z) = \sum_{(x,y) \,:\, g(x,y) = z} p_{X,Y}(x, y) \]
\end{formula}

\begin{esempio}{Due coppie e la stessa formula}
Si considerino due variabili doppie, $(X_1, Y_1) \in \{0, 1\}^2$ e $(X_2, Y_2) \in \{-1, 1\}^2$, con le pmf congiunte
riportate qui sotto (la prima è quella del capitolo precedente), e le due variabili
\[ Z_1 = 3X_1^2 + Y_1 , \qquad\qquad Z_2 = 3X_2^2 + Y_2 . \]
\begin{center}
\begin{tikzpicture}
  \begin{scope}[x=2.3cm, y=1.7cm]
    \draw[-{Stealth}] (-0.4,0) -- (1.5,0) node[below, font=\small] {$x_1$};
    \draw[-{Stealth}] (0,-0.4) -- (0,1.5) node[left, font=\small] {$y_1$};
    \foreach \x/\y/\p/\z in {0/0/\frac13/0, 0/1/\frac29/1, 1/0/\frac19/3, 1/1/\frac13/4} {
      \fill[mspiA] (\x,\y) circle (2.2pt);
      \node[font=\scriptsize, anchor=south west, align=left] at (\x,\y) {$p = \p$\\ {\color{mspiB!85!black}$z_1 = \z$}};
    }
    \node[font=\scriptsize, below] at (1,0) {1}; \node[font=\scriptsize, left] at (0,1) {1};
    \node[font=\small] at (0.6,-0.75) {$Z_1 = 3X_1^2 + Y_1$: quattro valori diversi};
  \end{scope}
  \begin{scope}[shift={(8.2,0.85)}, x=1.5cm, y=1.1cm]
    \draw[-{Stealth}] (-1.6,0) -- (1.8,0) node[below, font=\small] {$x_2$};
    \draw[-{Stealth}] (0,-1.5) -- (0,1.6) node[left, font=\small] {$y_2$};
    \draw[mspiB, rounded corners=6pt, dashed] (-1.25,0.72) rectangle (1.5,1.32);
    \draw[mspiC!75!black, rounded corners=6pt, dashed] (-1.25,-1.32) rectangle (1.5,-0.72);
    \foreach \x/\y/\p in {-1/-1/\frac14, -1/1/\frac12, 1/-1/\frac18, 1/1/\frac18} {
      \fill[mspiA] (\x,\y) circle (2.2pt);
      \node[font=\scriptsize, anchor=west] at (\x+0.06,\y) {$\p$};
    }
    \node[font=\scriptsize, mspiB!85!black, anchor=west] at (1.55,1.02) {$z_2 = 4$};
    \node[font=\scriptsize, mspiC!75!black, anchor=west] at (1.55,-1.02) {$z_2 = 2$};
    \node[font=\small] at (0.2,-1.95) {$Z_2 = 3X_2^2 + Y_2$: due valori soli};
  \end{scope}
\end{tikzpicture}
\end{center}
\textbf{Prima coppia.} I valori possibili di $Z_1$ sono $3 \cdot 0 + 0 = 0$, $3 \cdot 0 + 1 = 1$, $3 \cdot 1 + 0 = 3$,
$3 \cdot 1 + 1 = 4$: quindi $\card{\mathcal{Z}_1} = \card{\{0, 1, 3, 4\}} = 4 = \card{\{0,1\}}^2 = \card{\mathcal{X}}
\cdot \card{\mathcal{Y}}$, e la corrispondenza è biunivoca:
\[
  p_{Z_1}(0) = p_{X_1,Y_1}(0, 0) = \tfrac13 , \quad
  p_{Z_1}(1) = p_{X_1,Y_1}(0, 1) = \tfrac29 , \quad
  p_{Z_1}(3) = p_{X_1,Y_1}(1, 0) = \tfrac19 , \quad
  p_{Z_1}(4) = p_{X_1,Y_1}(1, 1) = \tfrac13 .
\]
\textbf{Seconda coppia.} Qui $X_2 = \pm 1$, quindi $X_2^2 = 1$ sempre e $Z_2 = 3 + Y_2$: i valori possibili sono solo
$\{2, 4\}$, e $\card{\mathcal{Z}_2} = 2 < \card{\{-1, 1\}}^2 = 4$. La corrispondenza non è biunivoca e le probabilità
collassano:
\[
  p_{Z_2}(2) = p_{X_2,Y_2}(1, -1) + p_{X_2,Y_2}(-1, -1) = \tfrac18 + \tfrac14 = \tfrac38 , \qquad
  p_{Z_2}(4) = p_{X_2,Y_2}(-1, 1) + p_{X_2,Y_2}(1, 1) = \tfrac12 + \tfrac18 = \tfrac58 .
\]
\end{esempio}

% ============================================================
\section{La media di una funzione di due variabili}
% ============================================================

Con riferimento alla trasformazione generica $Z = g(X, Y)$, il teorema fondamentale per il calcolo della media
(Teorema~\ref{teo:valoremedio}) si traduce in modo immediato.

\teorema[valoremedio2]{Teorema fondamentale per due variabili}{
  \[ E[Z] = E[g(X, Y)] = \sum_{x \in \mathcal{X}} \sum_{y \in \mathcal{Y}} g(x, y)\, p_{X,Y}(x, y)
          = \sum_{(x, y) \in \mathcal{X} \times \mathcal{Y}} g(x, y)\, p_{X,Y}(x, y) . \]
}

\begin{dimostrazione}
\passo{} È la stessa dimostrazione del caso di una variabile, con le coppie $(x, y)$ al posto dei valori $x$. Si
scrive $E[Z] = \sum_z z\, p_Z(z) = \sum_z \sum_{(x,y) \in A(z)} g(x, y)\, p_{X,Y}(x, y)$; gli insiemi $A(z)$ dividono
$\mathcal{X} \times \mathcal{Y}$ in gruppi disgiunti che lo coprono tutto, quindi sommare gruppo per gruppo equivale a
sommare su tutte le coppie.
\end{dimostrazione}

La conseguenza più importante riguarda le combinazioni lineari.

\teorema[linearita]{Linearità della media}{
  Se $a$ e $b$ sono costanti, per qualsiasi coppia di variabili aleatorie $X$ e $Y$:
  \[ E[aX + bY] = a\,E[X] + b\,E[Y] . \]
  Più in generale, se $X_1, X_2, \ldots, X_m$ sono $m$ variabili aleatorie con pmf congiunta
  $p_{X_1, \ldots, X_m}(x_1, \ldots, x_m)$ e $a_1, \ldots, a_m$ sono costanti:
  \[ E\left[ \sum_{i=1}^{m} a_i X_i \right] = \sum_{i=1}^{m} a_i\, E[X_i] . \]
}

\begin{dimostrazione}
\passo{Il teorema fondamentale con $g(x, y) = ax + by$}
\[ E[aX + bY] = \sum_{x \in \mathcal{X}} \sum_{y \in \mathcal{Y}} (ax + by)\, p_{X,Y}(x, y)
   = a \sum_{x} \sum_{y} x\, p_{X,Y}(x, y) + b \sum_{x} \sum_{y} y\, p_{X,Y}(x, y) . \]
\passo{Si raccoglie e si marginalizza} Nella prima somma $x$ non dipende da $y$ e si porta fuori dalla somma interna;
nella seconda si scambia l'ordine delle somme e si porta fuori $y$:
\[ E[aX + bY] = a \sum_{x \in \mathcal{X}} x \underbrace{\sum_{y \in \mathcal{Y}} p_{X,Y}(x, y)}_{p_X(x)}
   + b \sum_{y \in \mathcal{Y}} y \underbrace{\sum_{x \in \mathcal{X}} p_{X,Y}(x, y)}_{p_Y(y)} . \]
Le somme interne sono le marginali (Teorema~\ref{teo:marginali}).
\passo{Si riconoscono le medie}
\[ E[aX + bY] = a \sum_{x} x\, p_X(x) + b \sum_{y} y\, p_Y(y) = a\,E[X] + b\,E[Y] . \]
\passo{Con $m$ variabili} Si applica il risultato ripetutamente (è un'induzione): $\sum_{i=1}^{m} a_i X_i$ è la somma
di $\sum_{i=1}^{m-1} a_i X_i$ e di $a_m X_m$, e la media di una somma di due variabili è la somma delle medie.
\end{dimostrazione}

\warning{La linearità non chiede l'indipendenza}{
  La media di una somma è la somma delle medie \textbf{sempre}, qualunque sia la pmf congiunta: le variabili possono
  essere dipendenti quanto si vuole. È questo che la rende così utile. (Vedremo invece che per la varianza di una somma
  la dipendenza conta, eccome.)
}

\begin{esempio}{La media della binomiale, in una riga}
Nel Capitolo~5 abbiamo dimostrato $E[X_N] = Np$ con una lunga catena di passaggi. Con la linearità basta osservare che
il numero di teste in $N$ lanci è la somma di $N$ variabili di Bernoulli:
\[ X_N = B_1 + B_2 + \cdots + B_N , \qquad B_i = \begin{cases} 1 & \text{se l'$i$-esimo lancio dà testa} \\
   0 & \text{altrimenti} \end{cases} \]
Ogni $B_i$ vale 1 con probabilità $p$, quindi $E[B_i] = p$, e
\[ E[X_N] = E[B_1] + E[B_2] + \cdots + E[B_N] = N p . \]
\end{esempio}

\begin{esempio}{Le medie di $Z_1$ e $Z_2$}
Con le pmf trovate nell'esempio precedente:
\[
  E[Z_1] = 0 \cdot \tfrac13 + 1 \cdot \tfrac29 + 3 \cdot \tfrac19 + 4 \cdot \tfrac13 = \frac{2 + 3 + 12}{9} =
  \frac{17}{9} , \qquad
  E[Z_2] = 2 \cdot \tfrac38 + 4 \cdot \tfrac58 = \frac{6 + 20}{8} = \frac{26}{8} = \frac{13}{4} .
\]
\end{esempio}

% ============================================================
\section{Il teorema della media condizionata}
% ============================================================

Il teorema fondamentale si può riorganizzare in un modo molto utile. Con riferimento a $Z = g(X, Y)$, scriviamo la
congiunta come condizionata per marginale, $p_{X,Y}(x, y) = p_{X \mid Y}(x \mid y)\, p_Y(y)$:
\[
  E[Z] = \sum_{x \in \mathcal{X}} \sum_{y \in \mathcal{Y}} g(x, y)\, p_{X \mid Y}(x \mid y)\, p_Y(y)
  = \sum_{y \in \mathcal{Y}} p_Y(y) \underbrace{\sum_{x \in \mathcal{X}} g(x, y)\, p_{X \mid Y}(x \mid y)}_{h(y)}
  = \sum_{y \in \mathcal{Y}} p_Y(y)\, h(y) ,
\]
dove
\[ h(y) = \sum_{x \in \mathcal{X}} g(x, y)\, p_{X \mid Y}(x \mid y) = E[g(X, Y) \mid Y = y] \]
è la media di $g(X, Y)$ eseguita rispetto alla pmf condizionata $p_{X \mid Y}(x \mid y)$: la media di $g$ sapendo che
$Y = y$. Al variare di $y$ è una funzione di $y$, quindi $h(Y) = E[g(X, Y) \mid Y]$ è a sua volta una variabile
aleatoria (una funzione di $Y$), e l'ultima somma è la sua media.

\teorema[mediacondizionata]{Teorema della media condizionata}{
  \[ E[g(X, Y)] = E[h(Y)] = E\big[\, E[g(X, Y) \mid Y] \,\big] . \]
  La media si può calcolare in due tempi: prima la media condizionata a ogni valore di $Y$, poi la media di questi
  numeri, pesati con $p_Y$. Ovviamente i ruoli di $X$ e $Y$ si possono scambiare.
}

È lo stesso principio del Teorema~\ref{teo:mediacond} (la media è la media pesata delle medie condizionate), con la
partizione $\Omega = \bigcup_y \{Y = y\}$.

\begin{esempio}{La media di $Z_1$, in due tempi}
Calcoliamo di nuovo $E[Z_1] = E[E[Z_1 \mid Y_1]]$, dove $Z_1 = g(X_1, Y_1) = 3X_1^2 + Y_1$. Serve la pmf condizionata
$p_{X_1 \mid Y_1}(x_1 \mid y_1) = p_{X_1,Y_1}(x_1, y_1)/p_{Y_1}(y_1)$, e per questa la marginale di $Y_1$:
\[
  p_{Y_1}(0) = p_{X_1,Y_1}(0, 0) + p_{X_1,Y_1}(1, 0) = \tfrac13 + \tfrac19 = \tfrac49 , \qquad
  p_{Y_1}(1) = p_{X_1,Y_1}(0, 1) + p_{X_1,Y_1}(1, 1) = \tfrac29 + \tfrac13 = \tfrac59 .
\]
\begin{center}
\renewcommand{\arraystretch}{1.5}
\begin{tabular}{@{}cc@{}}
  \toprule
  $x_1 y_1$ & $p_{X_1 \mid Y_1}(x_1 \mid y_1) = \dfrac{p_{X_1,Y_1}(x_1, y_1)}{p_{Y_1}(y_1)}$ \\
  \midrule
  00 & $\frac13 \cdot \frac94 = \frac34$ \\
  01 & $\frac29 \cdot \frac95 = \frac25$ \\
  10 & $\frac19 \cdot \frac94 = \frac14$ \\
  11 & $\frac13 \cdot \frac95 = \frac35$ \\
  \bottomrule
\end{tabular}
\end{center}
Ora le medie condizionate. Sapendo che $Y_1 = 0$, la $Y_1$ nella formula vale 0; e siccome $X_1$ vale 0 o 1, il termine
$3X_1^2$ vale 3 solo quando $X_1 = 1$:
\[
  E[Z_1 \mid Y_1 = 0] = E[3X_1^2 + Y_1 \mid Y_1 = 0] = E[3X_1^2 \mid Y_1 = 0] = 3 \cdot 1^2 \cdot p_{X_1 \mid Y_1}(1 \mid 0)
  = 3 \cdot \tfrac14 = \tfrac34 .
\]
Sapendo che $Y_1 = 1$, invece, la $Y_1$ vale 1:
\[
  E[Z_1 \mid Y_1 = 1] = E[3X_1^2 + 1 \mid Y_1 = 1] = 3 \cdot 1^2 \cdot p_{X_1 \mid Y_1}(1 \mid 1) + 1 = 3 \cdot \tfrac35
  + 1 = \tfrac95 + 1 = \tfrac{14}{5} .
\]
Infine la media delle medie condizionate:
\[
  E[Z_1] = E[Z_1 \mid Y_1 = 0]\, p_{Y_1}(0) + E[Z_1 \mid Y_1 = 1]\, p_{Y_1}(1) = \tfrac34 \cdot \tfrac49 + \tfrac{14}{5}
  \cdot \tfrac59 = \tfrac39 + \tfrac{14}{9} = \tfrac{17}{9} ,
\]
lo stesso valore trovato prima.
\end{esempio}

% ============================================================
\section{Correlazione e covarianza}
% ============================================================

Le coppie $(\mu_X, \sigma_X^2)$ e $(\mu_Y, \sigma_Y^2)$ contengono informazioni globali, anche se sommarie, sulle due
pmf marginali. L'equivalente per la pmf congiunta è la \textbf{covarianza}, che dà un'idea, ancora abbastanza
sommaria, del grado di «dipendenza» fra $X$ e $Y$.

\dfn{Correlazione e covarianza}{
  La \term{correlazione} fra $X$ e $Y$ è la media del loro prodotto:
  \[ R_{XY} = E[XY] = \sum_{x \in \mathcal{X}} \sum_{y \in \mathcal{Y}} x\,y\, p_{X,Y}(x, y) . \]
  La \term{covarianza} fra $X$ e $Y$ è la media del prodotto degli scarti dalle rispettive medie:
  \[ \sigma_{XY}^2 = COV(X, Y) = E\left[(X - \mu_X)(Y - \mu_Y)\right]
     = \sum_{x} \sum_{y} (x - \mu_X)(y - \mu_Y)\, p_{X,Y}(x, y) . \]
}

\begin{figure}[H]
\centering
\begin{tikzpicture}
  \fill[mspiC!12] (0,0) rectangle (3,2.2); \fill[mspiC!12] (-3,-2.2) rectangle (0,0);
  \fill[mspiB!12] (-3,0) rectangle (0,2.2); \fill[mspiB!12] (0,-2.2) rectangle (3,0);
  \draw[-{Stealth}, mspigray] (-3.2,0) -- (3.3,0) node[below, font=\small, black] {$x$};
  \draw[-{Stealth}, mspigray] (0,-2.4) -- (0,2.5) node[left, font=\small, black] {$y$};
  \node[font=\scriptsize, below left] at (0,0) {$(\mu_X, \mu_Y)$};
  \node[font=\small, mspiC!60!black, align=center] at (1.5,1.6) {$x > \mu_X$, $y > \mu_Y$\\ prodotto $> 0$};
  \node[font=\small, mspiC!60!black, align=center] at (-1.5,-1.6) {$x < \mu_X$, $y < \mu_Y$\\ prodotto $> 0$};
  \node[font=\small, mspiB!80!black, align=center] at (-1.5,1.6) {$x < \mu_X$, $y > \mu_Y$\\ prodotto $< 0$};
  \node[font=\small, mspiB!80!black, align=center] at (1.5,-1.6) {$x > \mu_X$, $y < \mu_Y$\\ prodotto $< 0$};
  \node[font=\small, anchor=west, text width=6.4cm] at (3.7,0)
    {Gli assi sono centrati sulle medie. Il prodotto degli scarti $(x - \mu_X)(y - \mu_Y)$ è positivo nei quadranti
     verdi (le due variabili sono entrambe sopra o entrambe sotto la media) e negativo in quelli arancioni (una sopra,
     l'altra sotto). La covarianza è la media di questo prodotto: è positiva se la probabilità sta soprattutto nei
     quadranti verdi, negativa se sta soprattutto in quelli arancioni.};
\end{tikzpicture}
\caption{Il segno della covarianza.}
\end{figure}

\teorema[covarianza]{Covarianza e correlazione}{
  \[ COV(X, Y) = R_{XY} - \mu_X\, \mu_Y = E[XY] - E[X]\,E[Y] . \]
  In particolare, se almeno una delle due variabili ha media nulla, $COV(X, Y) = R_{XY}$.
}

\begin{dimostrazione}
\passo{Si sviluppa il prodotto} $(X - \mu_X)(Y - \mu_Y) = XY - \mu_X Y - \mu_Y X + \mu_X \mu_Y$.
\passo{Linearità della media} Le medie $\mu_X$ e $\mu_Y$ sono costanti, quindi
\[ COV(X, Y) = E[XY] - \mu_X\, E[Y] - \mu_Y\, E[X] + \mu_X \mu_Y = R_{XY} - \mu_X \mu_Y - \mu_X \mu_Y + \mu_X \mu_Y
   = R_{XY} - \mu_X \mu_Y . \]
\end{dimostrazione}

\dfn{Variabili incorrelate}{
  Due variabili aleatorie che abbiano covarianza nulla, $COV(X, Y) = 0$, si dicono \term{incorrelate}.
}

\teorema[indipincorr]{Indipendenti implica incorrelate}{
  Se $X$ e $Y$ sono indipendenti, allora sono incorrelate: $COV(X, Y) = 0$, e quindi $E[XY] = E[X]\,E[Y]$.
}

\begin{dimostrazione}
\passo{La congiunta si fattorizza} Se $p_{X,Y}(x, y) = p_X(x)\, p_Y(y)$:
\[ COV(X, Y) = \sum_{x} \sum_{y} (x - \mu_X)(y - \mu_Y)\, p_X(x)\, p_Y(y) . \]
\passo{La somma doppia si spezza} Il termine si fattorizza in una parte che dipende solo da $x$ e una che dipende solo
da $y$, quindi la somma doppia è il prodotto di due somme:
\[ COV(X, Y) = \sum_{x} (x - \mu_X)\, p_X(x) \sum_{y} (y - \mu_Y)\, p_Y(y) = E[X - \mu_X]\; E[Y - \mu_Y] = 0 \cdot 0 = 0 , \]
perché la media dello scarto dalla media è sempre nulla.
\end{dimostrazione}

La proposizione \textbf{non si inverte}: l'incorrelazione non implica, in generale, l'indipendenza.

\begin{esempio}{Incorrelate ma dipendenti}
Sia $X$ uniforme su $\{-1, 0, 1\}$ e sia $Y = X^2$. Le due variabili sono fortissimamente dipendenti: conoscendo $X$ si
conosce $Y$ con certezza (per esempio $P(Y = 0 \mid X = 0) = 1$, mentre $P(Y = 0) = \frac13$). Eppure
\[ E[X] = \tfrac13(-1 + 0 + 1) = 0 , \qquad E[XY] = E[X^3] = \tfrac13\left((-1)^3 + 0^3 + 1^3\right) = 0 , \]
quindi $COV(X, Y) = E[XY] - E[X]\,E[Y] = 0 - 0 = 0$: sono incorrelate. La covarianza vede solo la tendenza a crescere
o calare \emph{insieme}; qui $Y$ è grande sia quando $X$ è molto negativa sia quando è molto positiva, e i due effetti
si compensano.
\end{esempio}

% ============================================================
\section{Il coefficiente di correlazione}
% ============================================================

La covarianza ha un difetto: dipende dalle unità di misura (se $X$ si misura in centimetri invece che in metri, la
covarianza diventa cento volte più grande). Per avere un numero «puro» la si divide per le deviazioni standard; il
risultato sta sempre fra $-1$ e $1$.

\teorema[cauchy]{Limite della covarianza}{
  Per due variabili aleatorie qualsiasi (con $\sigma_X, \sigma_Y > 0$):
  \[ \abs{COV(X, Y)} \leq \sigma_X\, \sigma_Y . \]
}

\begin{dimostrazione}
\passo{Un quadrato ha media non negativa} Per entrambe le scelte del segno, la variabile
$\left(\frac{X - \mu_X}{\sigma_X} \pm \frac{Y - \mu_Y}{\sigma_Y}\right)^2$ è non negativa, quindi lo è la sua media:
\[ 0 \leq E\left[ \left( \frac{X - \mu_X}{\sigma_X} \pm \frac{Y - \mu_Y}{\sigma_Y} \right)^2 \right] . \]
\passo{Si sviluppa il quadrato} Con $(a \pm b)^2 = a^2 + b^2 \pm 2ab$ e la linearità della media:
\[ 0 \leq E\left[ \left( \frac{X - \mu_X}{\sigma_X} \right)^2 \right] + E\left[ \left( \frac{Y - \mu_Y}{\sigma_Y}
   \right)^2 \right] \pm 2\, E\left[ \frac{X - \mu_X}{\sigma_X} \cdot \frac{Y - \mu_Y}{\sigma_Y} \right] . \]
\passo{Si riconoscono i pezzi} Il primo termine è $\frac{E[(X - \mu_X)^2]}{\sigma_X^2} = \frac{\sigma_X^2}{\sigma_X^2} =
1$, il secondo allo stesso modo vale 1, il terzo è $\pm 2\,\frac{COV(X, Y)}{\sigma_X \sigma_Y}$:
\[ 0 \leq 1 + 1 \pm 2\, \frac{COV(X, Y)}{\sigma_X\, \sigma_Y} \quad\iff\quad -2 \leq \pm 2\, \frac{COV(X, Y)}{\sigma_X\,
   \sigma_Y} \quad\iff\quad -\sigma_X \sigma_Y \leq \pm\, COV(X, Y) . \]
\passo{I due segni insieme} Con il segno $+$: $COV(X, Y) \geq -\sigma_X \sigma_Y$. Con il segno $-$:
$-COV(X, Y) \geq -\sigma_X \sigma_Y$, cioè $COV(X, Y) \leq \sigma_X \sigma_Y$. Insieme:
\[ -\sigma_X \sigma_Y \leq COV(X, Y) \leq \sigma_X \sigma_Y \quad\iff\quad \abs{COV(X, Y)} \leq \sigma_X\, \sigma_Y . \]
\end{dimostrazione}

\dfn{Coefficiente di correlazione}{
  Il \term{coefficiente di correlazione} fra $X$ e $Y$ è
  \[ \rho_{XY} = \frac{COV(X, Y)}{\sigma_X\, \sigma_Y} \in [-1, 1] . \]
}

\begin{figure}[H]
\centering
\begin{tikzpicture}
\begin{groupplot}[group style={group size=4 by 1, horizontal sep=0.35cm},
  width=4.4cm, height=4.4cm, xmin=-2.7, xmax=2.7, ymin=-2.7, ymax=2.7, xtick=\empty, ytick=\empty,
  axis lines=middle, axis line style={mspigray}, title style={font=\small}]
\nextgroupplot[title={$\rho \approx 0{,}9$}]
  \addplot[only marks, mark=*, mark size=0.9pt, mspiA] coordinates {(-1.45,-1.20) (-0.28,-0.46) (0.97,0.90) (-0.27,0.54) (-0.49,-1.00) (-1.37,-1.47) (0.66,0.96) (1.14,1.83) (1.72,1.84) (1.36,1.52) (-0.53,-0.10) (0.11,0.03) (2.06,2.22) (-0.44,-0.40) (2.41,2.09) (0.80,1.01) (0.57,0.80) (-0.43,-0.57) (-0.11,-0.23) (0.87,1.49) (-0.84,-0.49) (-0.94,-1.13) (-0.19,-0.10) (-0.39,-1.48) (1.46,1.22) (-0.59,-0.20) (-0.55,-1.10) (0.55,0.71) (-0.78,-0.31) (0.43,-0.09) (1.83,1.79) (-1.10,-0.76) (0.18,0.62) (-1.31,-0.70) (1.01,1.24) (0.67,1.21) (0.78,0.79) (-0.79,-0.20) (1.05,1.10) (0.21,0.53) (-2.50,-2.08) (-2.48,-2.10) (1.12,1.87) (-0.43,-0.17) (-0.02,-0.98) (0.79,0.73) (0.12,-0.29) (0.01,0.06) (1.35,1.36) (-0.22,-0.20) (-0.77,-0.01) (0.70,0.09) (0.16,0.67) (-1.39,-1.16) (-0.19,0.88) (0.47,0.26) (-0.52,0.39) (-0.42,0.65) (-1.54,-1.45) (-0.19,-0.47) (-0.17,0.18) (-0.63,-1.18) (1.41,1.94) (-0.65,-0.43) (-2.56,-1.82) (-1.30,-1.40) (2.24,1.48) (0.07,-0.39) (-1.83,-2.14) (0.32,-0.27) (-0.55,-0.15) (-0.11,-0.66) (1.56,1.48) (1.47,0.88) (1.76,1.01) (-1.72,-1.72) (-0.33,-0.16) (-0.52,-0.17) (1.40,0.99) (-1.35,-1.37) (0.88,1.17) (-0.31,-0.69) (0.03,0.47) (-0.68,-0.02) (0.27,0.39) (-0.59,0.09) (0.09,-0.56) (0.87,0.62) (1.14,1.12) (-0.59,-0.44) (-0.45,-0.75) (0.07,0.45) (0.12,0.09) (-0.52,-0.19) (0.32,0.73) (0.41,1.33) (-1.55,-1.40) (0.58,0.94) (-1.51,-1.88) (0.00,0.48) (1.16,0.39) (0.18,0.69) (0.07,-0.37) (-0.68,-1.18) (1.25,1.03) (-0.39,-0.72) (-0.10,-0.47) (-0.95,-0.85) (-0.19,0.05) (0.19,0.71)};
\nextgroupplot[title={$\rho \approx 0{,}3$}]
  \addplot[only marks, mark=*, mark size=0.9pt, mspiA] coordinates {(-0.00,0.29) (0.27,-0.63) (1.24,-0.06) (-1.56,0.68) (0.70,-0.25) (1.01,1.31) (0.75,-0.81) (0.22,-0.87) (-0.58,-0.70) (-0.81,0.22) (-0.76,0.36) (0.21,-0.46) (0.21,0.76) (1.20,0.21) (1.80,1.21) (1.30,-0.79) (-1.03,0.82) (-0.19,0.41) (1.25,-0.89) (0.50,2.05) (0.85,0.42) (0.61,0.26) (0.82,1.56) (-0.19,0.52) (0.35,0.86) (0.89,-0.67) (1.22,1.28) (-1.06,-0.60) (-1.43,-0.74) (0.37,-0.93) (-0.74,-1.05) (-0.66,-0.92) (-1.04,0.21) (-0.18,-0.80) (1.04,-0.13) (0.94,-0.80) (1.46,0.13) (-0.91,0.10) (-1.31,-1.44) (0.28,-0.32) (0.44,0.33) (0.29,0.26) (0.76,0.30) (1.15,1.33) (-1.94,-2.30) (1.61,0.18) (0.23,0.01) (0.32,2.23) (-0.67,-0.14) (0.04,-0.97) (-0.84,-0.51) (0.02,0.41) (-1.93,-2.12) (1.70,0.46) (0.87,-0.06) (-0.68,1.79) (-1.20,-1.24) (0.72,0.30) (-1.61,-0.40) (0.77,-1.13) (-0.52,0.59) (1.96,-1.47) (-0.66,-0.17) (-0.88,-1.13) (-0.47,0.35) (-1.03,0.49) (0.87,1.14) (1.58,0.10) (-0.01,0.97) (1.31,1.38) (-0.42,-0.57) (-1.29,-0.82) (-0.86,-0.73) (-0.34,0.70) (0.90,0.30) (-1.11,-0.40) (0.47,1.00) (0.25,0.88) (-0.62,0.43) (0.78,-0.03) (-0.72,0.94) (-0.96,0.34) (-0.23,-0.38) (-0.22,-0.46) (-0.00,1.71) (1.49,0.43) (0.51,1.06) (-0.39,1.62) (0.07,0.64) (-0.92,-1.34) (-0.96,-0.91) (1.07,0.11) (0.92,-0.04) (-0.92,-0.13) (-0.36,-1.65) (0.08,1.17) (-0.17,-1.15) (-0.72,-1.36) (-0.21,-0.93) (-0.08,0.17) (1.36,-0.57) (-2.02,0.85) (0.47,-0.14) (0.54,-0.65) (-0.30,0.56) (-0.28,0.06) (0.47,0.31) (-0.21,0.44) (1.36,1.70) (-0.84,0.45)};
\nextgroupplot[title={$\rho \approx 0$}]
  \addplot[only marks, mark=*, mark size=0.9pt, mspiA] coordinates {(-1.87,-1.38) (0.44,0.57) (1.40,-1.69) (1.12,-0.51) (-0.43,0.48) (1.44,-0.91) (-1.68,-0.56) (-0.08,1.23) (0.40,-0.27) (-0.47,-0.43) (1.00,-1.63) (-0.66,-0.84) (1.16,0.16) (-0.64,0.78) (1.99,1.05) (-1.27,-0.18) (1.17,-1.78) (-1.19,-0.38) (0.55,-0.79) (0.18,-0.13) (0.35,-1.17) (-2.19,-2.12) (0.19,1.00) (0.51,-1.58) (-1.71,-0.26) (-0.89,-1.81) (-0.56,-0.77) (-0.09,0.40) (0.21,0.21) (0.54,-0.08) (-0.43,-0.32) (1.21,-0.26) (0.83,1.68) (1.76,-2.13) (0.17,1.00) (-0.43,-0.69) (-1.41,-0.66) (0.02,-1.86) (0.00,-0.82) (0.04,-1.07) (0.10,0.18) (-0.68,-0.65) (0.96,0.65) (-0.76,0.45) (-1.14,-0.50) (0.70,-0.47) (0.60,-0.63) (0.32,-0.13) (-0.96,2.11) (-0.83,-1.33) (0.14,0.49) (0.36,-0.09) (0.48,0.95) (0.64,-1.37) (1.74,-0.02) (-2.03,0.85) (-0.13,0.30) (2.00,-0.52) (-0.61,-0.97) (0.20,0.67) (0.40,0.18) (-2.11,0.28) (1.09,-0.50) (-0.52,-1.00) (0.26,-1.54) (1.92,-0.18) (-0.58,-1.20) (0.55,0.99) (-0.55,0.19) (0.31,0.98) (0.63,1.25) (-0.49,-1.10) (-0.91,0.40) (0.79,2.44) (1.24,-0.75) (-0.33,-0.45) (-0.58,-0.16) (-0.60,1.19) (0.35,0.64) (-0.14,-1.14) (0.45,-0.38) (2.01,1.25) (-1.19,-0.97) (-0.53,-0.20) (1.16,-0.16) (0.47,-0.30) (-0.20,1.25) (0.03,1.98) (-0.03,0.40) (0.41,0.68) (1.94,-0.66) (-0.51,-0.92) (1.82,0.49) (0.47,0.48) (-1.83,0.79) (0.15,0.33) (0.11,-0.07) (-0.05,-0.55) (-0.21,0.65) (-2.09,0.84) (0.66,-1.05) (0.53,0.33) (-0.50,1.71) (0.20,-0.11) (-1.28,2.19) (0.24,0.33) (-0.21,2.09) (-0.13,1.59) (1.07,1.00) (0.65,1.17)};
\nextgroupplot[title={$\rho \approx -0{,}6$}]
  \addplot[only marks, mark=*, mark size=0.9pt, mspiB] coordinates {(-0.10,-1.40) (-2.03,2.39) (0.63,-0.09) (-0.69,-0.92) (-0.07,1.27) (-0.09,0.07) (0.10,-0.03) (0.31,1.43) (-0.86,-0.69) (0.86,-0.96) (-0.48,-0.64) (-0.16,-0.18) (0.45,-0.58) (-1.44,1.37) (-0.06,-1.29) (-0.32,0.78) (-0.04,-0.19) (0.61,-0.87) (-0.68,0.86) (-0.14,-0.33) (0.24,0.03) (0.95,0.07) (0.11,-0.08) (-0.91,0.79) (-0.34,0.85) (0.82,-0.10) (-0.04,0.27) (1.52,0.16) (-0.57,-1.60) (0.64,-1.52) (0.07,0.29) (1.51,-1.17) (-0.09,-0.33) (0.32,-1.08) (0.59,-0.15) (1.37,-1.16) (0.35,-0.08) (0.00,-0.67) (-0.54,0.43) (0.51,0.40) (-0.58,0.60) (0.63,0.18) (0.27,-0.02) (-0.32,0.26) (-0.17,-1.10) (-0.68,-1.21) (0.53,0.27) (-0.33,0.37) (0.18,-0.29) (1.49,-1.16) (0.35,-0.27) (2.02,-2.56) (0.23,-1.09) (-0.68,-0.18) (0.79,-1.30) (-0.97,0.47) (-0.92,-0.04) (1.60,-1.57) (-0.77,1.54) (0.12,0.73) (-1.59,0.93) (0.51,-0.18) (0.57,0.51) (0.61,0.37) (0.78,-1.05) (0.22,0.79) (-0.91,1.04) (0.66,0.51) (-1.51,1.08) (2.27,-1.37) (1.94,-0.31) (0.19,-0.66) (-0.28,-0.70) (-1.34,1.24) (-1.41,1.27) (-0.04,-0.05) (-1.00,-0.15) (0.83,-1.24) (-0.54,0.28) (0.71,0.93) (1.99,-0.81) (0.70,0.00) (-0.18,0.00) (-1.69,1.60) (-0.46,0.62) (2.26,-0.95) (1.18,-0.18) (-0.53,-0.07) (-1.42,0.61) (0.20,0.38) (0.61,-0.65) (-1.10,0.83) (0.34,-1.63) (0.12,-1.08) (-0.42,0.22) (-0.95,0.32) (0.04,-0.72) (0.71,-0.01) (0.57,-1.13) (-0.17,-0.10) (0.14,0.45) (-0.50,0.39) (0.84,-0.81) (0.74,-0.52) (-1.38,0.96) (0.17,-0.40) (1.93,-1.61) (-1.44,0.91) (-0.20,0.52) (2.54,-1.66)};
\end{groupplot}
\end{tikzpicture}
\caption{Ogni punto è una coppia $(x, y)$ estratta da quattro coppie di variabili con coefficienti di correlazione
diversi (simulazione). Con $\rho$ vicino a 1 i punti si allineano lungo una retta crescente; con $\rho$ vicino a 0 non
c'è una tendenza; con $\rho$ negativo la tendenza è decrescente.}
\end{figure}

\info{Che cosa misura $\rho$}{
  Il coefficiente di correlazione misura la dipendenza \textbf{lineare}: quanto i punti $(x, y)$ tendono a disporsi
  lungo una retta. Dalla dimostrazione si vede che $\rho_{XY} = \pm 1$ esattamente quando la media del quadrato è zero,
  cioè quando $\frac{X - \mu_X}{\sigma_X} = \mp \frac{Y - \mu_Y}{\sigma_Y}$ con certezza: $Y$ è una funzione lineare di
  $X$, $Y = aX + b$. Un $\rho$ vicino a zero, invece, non esclude affatto una dipendenza non lineare, come
  nell'esempio $Y = X^2$.
}

% ============================================================
\section{La varianza di una combinazione lineare}
% ============================================================

\teorema[varsomma]{Varianza di una combinazione lineare}{
  Sia $Z = aX + bY$, con $a$ e $b$ costanti. Allora $\mu_Z = a\mu_X + b\mu_Y$ e
  \[ \sigma_Z^2 = a^2 \sigma_X^2 + b^2 \sigma_Y^2 + 2ab\, COV(X, Y) . \]
  In particolare, se $X$ e $Y$ sono incorrelate (per esempio indipendenti), $\sigma_{X+Y}^2 = \sigma_X^2 + \sigma_Y^2$.
}

\begin{dimostrazione}
\passo{Lo scarto di $Z$} La media segue dalla linearità; lo scarto dalla media è
\[ Z - \mu_Z = aX + bY - a\mu_X - b\mu_Y = a\,(X - \mu_X) + b\,(Y - \mu_Y) . \]
\passo{Il quadrato} Sviluppando $(a u + b v)^2 = a^2 u^2 + b^2 v^2 + 2ab\,uv$:
\[ \sigma_Z^2 = E\left[(Z - \mu_Z)^2\right] = E\left[ \left( a(X - \mu_X) + b(Y - \mu_Y) \right)^2 \right] \]
\[ = a^2\, E[(X - \mu_X)^2] + b^2\, E[(Y - \mu_Y)^2] + 2ab\, E[(X - \mu_X)(Y - \mu_Y)]
   = a^2 \sigma_X^2 + b^2 \sigma_Y^2 + 2ab\, COV(X, Y) . \]
\end{dimostrazione}

\begin{esempio}{La varianza della binomiale, in poche righe}
Riprendiamo $X_N = B_1 + \cdots + B_N$, con le $B_i$ di Bernoulli \textbf{indipendenti} (i lanci lo sono). Ogni $B_i$
ha varianza $p(1-p)$. Per due lanci, $\sigma_{B_1 + B_2}^2 = p(1-p) + p(1-p)$, perché la covarianza di due variabili
indipendenti è nulla. Aggiungendo un lancio alla volta: la somma dei primi $n - 1$ lanci è una funzione di quei lanci,
che sono indipendenti dall'$n$-esimo, e la covarianza di una somma si spezza nella somma delle covarianze (dalla
definizione e dalla linearità della media: $COV(U + V, W) = COV(U, W) + COV(V, W)$), tutte nulle. Così ogni lancio
aggiunge $p(1-p)$, e
\[ \sigma_{X_N}^2 = N\,p\,(1 - p) , \]
lo stesso risultato ottenuto nel Capitolo~7 con il trucco della derivata.
\end{esempio}

% ============================================================
\section{Quale coppia è più correlata?}
% ============================================================

\begin{esercizio}{Confrontare due coefficienti di correlazione}
Date le due pmf congiunte introdotte all'inizio del capitolo, per $(X_1, Y_1) \in \{0, 1\}^2$ e
$(X_2, Y_2) \in \{-1, 1\}^2$, trovare quale delle due dà luogo a un coefficiente di correlazione maggiore (in valore
assoluto).
\svolgimento
\textbf{Prima coppia.} Le marginali, sommando righe e colonne della tabella:
$p_{X_1}(0) = \frac59$, $p_{X_1}(1) = \frac49$; $p_{Y_1}(0) = \frac49$, $p_{Y_1}(1) = \frac59$. Quindi
\[
  E[X_1] = 0 \cdot \tfrac59 + 1 \cdot \tfrac49 = \tfrac49 , \qquad E[Y_1] = 0 \cdot \tfrac49 + 1 \cdot \tfrac59 = \tfrac59 ,
\]
\[
  E[X_1^2] = 0^2 \cdot \tfrac59 + 1^2 \cdot \tfrac49 = \tfrac49 , \qquad E[Y_1^2] = \tfrac59 ,
\]
\[
  \sigma_{X_1}^2 = E[X_1^2] - E^2[X_1] = \tfrac49 - \tfrac{16}{81} = \tfrac{36 - 16}{81} = \tfrac{20}{81} \approx 0{,}247 ,
  \qquad
  \sigma_{Y_1}^2 = \tfrac59 - \tfrac{25}{81} = \tfrac{45 - 25}{81} = \tfrac{20}{81} \approx 0{,}247 .
\]
Per $E[X_1 Y_1]$ si considera l'unico prodotto di valori diverso da zero, quello con $x_1 = y_1 = 1$:
\[ E[X_1 Y_1] = 1 \cdot 1 \cdot p_{X_1,Y_1}(1, 1) = \tfrac13 , \qquad
   COV(X_1, Y_1) = E[X_1 Y_1] - E[X_1]\,E[Y_1] = \tfrac13 - \tfrac49 \cdot \tfrac59 = \tfrac{27 - 20}{81} = \tfrac{7}{81}
   \approx 0{,}086 . \]
\[ \rho_{X_1 Y_1} = \frac{COV(X_1, Y_1)}{\sigma_{X_1} \sigma_{Y_1}} = \frac{7/81}{\sqrt{\frac{20}{81} \cdot
   \frac{20}{81}}} = \frac{7/81}{20/81} = \frac{7}{20} = 0{,}35 . \]
\textbf{Seconda coppia.} Le marginali: $p_{X_2}(-1) = \frac14 + \frac12 = \frac34$, $p_{X_2}(1) = \frac18 + \frac18 =
\frac14$; $p_{Y_2}(-1) = \frac14 + \frac18 = \frac38$, $p_{Y_2}(1) = \frac12 + \frac18 = \frac58$. Quindi
\[ E[X_2] = -1 \cdot \tfrac34 + 1 \cdot \tfrac14 = -\tfrac12 , \qquad E[Y_2] = -1 \cdot \tfrac38 + 1 \cdot \tfrac58 =
   \tfrac14 , \]
e, siccome $(\pm 1)^2 = 1$, $E[X_2^2] = E[Y_2^2] = 1$:
\[ \sigma_{X_2}^2 = 1 - \tfrac14 = \tfrac34 , \qquad \sigma_{Y_2}^2 = 1 - \tfrac1{16} = \tfrac{15}{16} . \]
Per $E[X_2 Y_2]$ servono tutti e quattro i punti (il prodotto vale $+1$ se i segni sono uguali, $-1$ se sono diversi):
\[
  E[X_2 Y_2] = (-1 \cdot 1)\,\tfrac12 + (-1 \cdot -1)\,\tfrac14 + (1 \cdot 1)\,\tfrac18 + (1 \cdot -1)\,\tfrac18
  = -\tfrac12 + \tfrac14 + \tfrac18 - \tfrac18 = -\tfrac14 ,
\]
\[ COV(X_2, Y_2) = -\tfrac14 - \left(-\tfrac12\right) \cdot \tfrac14 = -\tfrac14 + \tfrac18 = -\tfrac18 , \]
\[ \rho_{X_2 Y_2} = \frac{-1/8}{\sqrt{\frac34 \cdot \frac{15}{16}}} = \frac{-1/8}{\sqrt{45/64}} = \frac{-1/8}{\sqrt{45}/8}
   = -\frac{1}{\sqrt{45}} \approx -0{,}149 . \]
\textbf{Conclusione.} Essendo $\abs{\rho_{X_1 Y_1}} = 0{,}35 > \abs{\rho_{X_2 Y_2}} \approx 0{,}149$, la coppia
$(X_1, Y_1)$ risulta a correlazione assoluta maggiore. Il segno negativo della seconda ha un significato preciso:
ci si aspetta che $X_2$ e $Y_2$ tendano a discostarsi dalle rispettive medie in direzioni opposte (se una è superiore
alla media, l'altra tipicamente è inferiore, e viceversa).
\end{esercizio}

\gbox{In sintesi}{mspiA!80!black}{
\begin{itemize}
  \item Se $Z = g(X, Y)$: $p_Z(z) = \sum_{g(x,y) = z} p_{X,Y}(x, y)$ (le probabilità collassano se $g$ non è
        biunivoca sulle coppie), e $E[g(X, Y)] = \sum_x \sum_y g(x, y)\, p_{X,Y}(x, y)$.
  \item \textbf{Linearità}: $E[\sum a_i X_i] = \sum a_i E[X_i]$, sempre, anche per variabili dipendenti.
  \item \textbf{Media condizionata}: $E[g(X, Y)] = E\big[E[g(X, Y) \mid Y]\big]$: prima la media sapendo $Y$, poi la
        media su $Y$.
  \item Correlazione $R_{XY} = E[XY]$; \textbf{covarianza} $COV(X, Y) = E[(X - \mu_X)(Y - \mu_Y)] = R_{XY} -
        \mu_X \mu_Y$. Indipendenti $\Rightarrow$ incorrelate, non il viceversa.
  \item $\abs{COV(X, Y)} \leq \sigma_X \sigma_Y$; \textbf{coefficiente di correlazione}
        $\rho_{XY} = COV/(\sigma_X \sigma_Y) \in [-1, 1]$, misura della dipendenza lineare.
  \item $\sigma_{aX + bY}^2 = a^2 \sigma_X^2 + b^2 \sigma_Y^2 + 2ab\, COV(X, Y)$: per variabili incorrelate le varianze
        si sommano.
\end{itemize}
}
